\section{Conclusions and Outlook} \label{sec:conclusions}
We have described the Galactic Radio Explorer (GReX), a network of low-cost $1.4\,\text{GHz}$ radio antennas searching for bright transients between $32.768\,\mu\text{s}$ and $1.024\,\text{ms}$.
Terminals have been deployed at multiple sites around the world, though most are currently in the US.
Since the discovery of FRB\,200428, GReX and an upgraded STARE2 have more than quadrupled the total square degree hours devoted to ultra-bright FRB searching.
The improved system temperature of GReX LNAs over the original STARE2 hardware results in $\sim$\,twice the sensitivity, resulting in a deeper search.
We also search for extremely narrow bursts $\ll1\,\text{ms}$, which is not possible at lower frequencies or in systems with a large number of beams.
With no detections, we update the all-sky rate of FRBs above $\sim50\,\text{kJy}$ to be no larger than 8.5\,sky$^{-1}$\,yr$^{-1}$ at 95$\%$ confidence.
This indicates that FRB-like emission from Galactic sources is rare, and FRB\,200428 was even more unusual than initially thought. 

The GReX boxes have proven to be replicable and stable.
For example, the OVRO GReX terminal did not require on-site human intervention for more than a year of observation. Five sites are on-sky and continuously searching for bright FRBs.
These locations are  OVRO in California, the Cornell GReX in Ithaca, New York, Harvard GReX in Cambridge Massachusetts, Hat Creek GReX in California, and an Ireland station in Birr.
A box has been shipped and tesed in New South Wales, though not yet deployed.
In the future, we hope to expand to multiple sites in the Southern Hemisphere to improve exposure to the Galactic plane, where most magnetars reside. 
